\subsection{A predicted subtype--redshift march}
\label{sec:subtype_march}

The source-level closure of Section~\ref{sec:physical_visibility_closure}
assigns every synthetic source a dense photon fraction through the continuum
dominance $D$. Because $D$ is a lifecycle quantity---it grows as the nuclear
reservoir feeds while the stellar host is still assembling---the closure makes
a population-level prediction that the stacked comparison at $4.5\leq z\leq6.5$
does not use: the mixture of continuum subtypes must evolve with redshift.

We evaluate this prediction directly on the forward catalogue, weighting each
source by its selection probability and comoving weight. The predicted red
fraction (xLRD plus plusLRD) rises monotonically from $0.5\%$ at $z\simeq9.5$
to $10.0\%$ at $z\simeq7.5$, $15.2\%$ at $z\simeq5.5$, $38.8\%$ at
$z\simeq4.5$, and $47.8\%$ at $z\simeq3.25$, while the bLRD fraction falls
from $54\%$ to below $20\%$. The selection-weighted median
$f_{\mathrm{dense}}$ climbs from 0.40 to about 0.85 over the same interval.
Blue-continuum broad-line sources are therefore predicted to be the
\emph{majority} subtype at $z>8.5$, and xLRDs to be essentially absent there.

Two checks argue that this march is population physics rather than a
selection artifact. First, removing the selection weighting entirely
strengthens the trend: the intrinsic red fraction evolves from $0.4\%$ at
$z\simeq9.5$ to $70.5\%$ at $z\simeq3.25$, and every subtype trend retains its
sign. The gate sensitivity documented in
Section~\ref{sec:physical_visibility_closure} rescales the absolute
fractions, but a transparent gate cannot manufacture or erase a monotonic
factor-of-$\sim$100 evolution that survives with no gate at all. Second, the
rise from $z\simeq9.5$ to $z\simeq4.5$ is monotonic in all four random-seed
realizations; only the $z\simeq3.25$ bin, with an effective sample size near
17, shows appreciable seed scatter and should be read as indicative.

The march yields three falsifiable statements. (i) At $z>8.5$, bLRDs should
dominate the broad-line LRD population, and confirmed xLRDs should be rare at
the well-below-percent level; a modest confirmed xLRD sample at these
redshifts would contradict the closure. (ii) Between $z\simeq7.5$ and
$z\simeq4.5$ the red fraction should rise at roughly $0.08$--$0.10$ per unit
redshift, a slope measurable by splitting existing stacked and spectroscopic
samples at $z\simeq5.5$. (iii) Because the driver is $D$ rather than
luminosity, the $f_{\mathrm{dense}}$--redshift anticorrelation should persist
at fixed bolometric luminosity, distinguishing it from a pure
luminosity-selection trend.

Two caveats bound the claim. The $z>8.5$ bin lies in the regime where the
demographic model underpredicts the reported abundance by $2.4\sigma$; if a
missing early formation or luminosity channel operates there, it could alter
the high-redshift end of the march, and the $z\lesssim7$ portion of the
prediction is insulated from that tension. And the absolute subtype fractions
inherit the uncalibrated completeness mapping of
Section~\ref{sec:physical_visibility_closure}; the prediction defended here is
the ordered, monotonic evolution, not the individual percentages. The
calculation is reproduced by \texttt{analyze\_subtype\_march.py}, with
machine-readable output in \texttt{subtype\_march\_summary.json}.
